After the coherent analysis described above, the data are represented by a
collection of resampled spectra indexed by chunk time, by the coherent
resampling parameter $\beta$, and by Fourier frequency in the resampled
coordinate.  The purpose of the semicoherent stage is to combine these spectra
along the frequency evolution of a putative signal.  For a trial set of
long-timescale parameters $\theta$, the signal track predicts a frequency
$f_i(\theta)$ in the $i$th coherent chunk.  The stack-slide statistic is the
weighted sum of the coherent powers read from the corresponding bins,
\begin{equation}
    \Lambda(\theta)
    =
    \sum_{i\in{\cal I}(\theta)}
    w_i\,P_i\!\left[f_i(\theta)\right] .
    \label{eq:stack_slide_statistic}
\end{equation}
Here $P_i[f]$ denotes the coherent power in chunk $i$ at the nearest available
frequency bin, ${\cal I}(\theta)$ is the set of chunks for which the trial
track lies inside the analysis band, and $w_i$ denotes the semicoherent weight.
In the estimates below we take $w_i$ to be uniform.  In a search on real data
these weights should include gaps, detector noise, the NUFFT noise transfer
function, and the covariance of the five-vector sideband estimator.

The track parameters used for the current count estimate are
\begin{equation}
    \theta^a = \left(\ln M_c,\,t_{20}\right),
    \label{eq:track_coordinates}
\end{equation}
where $t_{20}$ is the time at which the 3.5PN TaylorF2 track crosses
$20\,{\rm Hz}$.  For a given chirp mass and crossing time, the track determines
the expected carrier frequency in every coherent chunk.  The Fourier frequency
axis produced by the NUFFT handles the carrier-frequency search within each
chunk, so there is no additional semicoherent $f_0$ dimension in this
parameterization.

\subsubsection{Mismatch for a power stack}
\label{sec:power_stack_mismatch}

A template bank is needed because the true values of $M_c$ and $t_{20}$ are not
known.  If a nearby template uses $\theta+\Delta\theta$ instead of $\theta$,
the track will read power from slightly displaced Fourier bins.  The natural
``mismatch'' for this stage is therefore not the coherent waveform mismatch,
but the fractional loss in the expected stack-slide power.

This is the main difference from the geometric placement formalism used in
standard compact-binary matched filtering.  In the usual construction, as in
the template-bank work of Brown et al. and the MANIFOLD work of Hanna et al.,
the local metric is induced by the noise-weighted waveform overlap,
\begin{equation}
    \mu_{\rm coh}
    =
    1
    -
    \max_{t_c,\phi_c}
    \left<
        \hat h(\theta)
        \middle|
        \hat h(\theta+\Delta\theta)
    \right> .
    \label{eq:coherent_waveform_mismatch}
\end{equation}
Our semicoherent statistic has already discarded the relative phase between
coherent chunks.  The geometry is therefore induced by a different map from
parameters to observables: the map from $\theta$ to the expected sequence of
coherent-bin powers along a track.  The two constructions are closely related
because both are local quadratic expansions of a loss function, but the loss
function is different.

Let
\begin{equation}
    \Delta f_i
    =
    \frac{\partial f_i}{\partial\theta^a}\Delta\theta^a
\end{equation}
be the track-frequency error in chunk $i$.  The quantity that controls the
power loss is this error in units of the local coherent bin width,
\begin{equation}
    d_i
    =
    \frac{\Delta f_i}{\Delta f_{{\rm bin},i}}
    =
    \frac{1}{\Delta f_{{\rm bin},i}}
    \frac{\partial f_i}{\partial\theta^a}\Delta\theta^a .
    \label{eq:semicoherent_bin_offset}
\end{equation}
For the coordinates in Eq.~\eqref{eq:track_coordinates},
\begin{equation}
    d_i
    =
    \frac{1}{\Delta f_{{\rm bin},i}}
    \frac{\partial f_i}{\partial \ln M_c}\Delta\ln M_c
    +
    \frac{1}{\Delta f_{{\rm bin},i}}
    \frac{\partial f_i}{\partial t_{20}}\Delta t_{20},
    \label{eq:semicoherent_bin_offset_coordinates}
\end{equation}
with $\partial f_i/\partial t_{20}=-\dot f_i$.

The coherent bin width is set by the duration of the chunk in the resampled
coordinate, not exactly by $T_{\rm coh}$.  For a 0PN template with chirp
parameter $\beta_i$ in chunk $i$,
\begin{equation}
    \Delta \tau_i
    =
    \frac{3}{5\beta_i}
    \left[
        1
        -
        \left(1-\frac{8}{3}\beta_i T_{\rm coh}\right)^{5/8}
    \right],
    \qquad
    \Delta f_{{\rm bin},i} = \frac{1}{\Delta\tau_i}.
    \label{eq:tau_bin_width}
\end{equation}

If a signal is offset from the nearest Fourier bin by $d$ bins, the recovered
power is reduced by a bin-response function $B(d)$.  For a rectangular coherent
window, $B(d)=\mathrm{sinc}^2(d)$.  Since the true signal frequency is not
generally known relative to the nearest-bin center, we also consider the
bin-averaged response
\begin{equation}
    B_{\rm avg}(d)
    =
    \int_{-1/2}^{1/2}
    \mathrm{sinc}^2(x+d)\,dx .
    \label{eq:bin_averaged_response}
\end{equation}
The expected retained stack-slide power is
\begin{equation}
    R(\theta,\Delta\theta)
    =
    \frac{
        \sum_{i\in{\cal I}(\theta)}
        w_i B(d_i)
    }{
        \sum_{i\in{\cal I}(\theta)}
        w_i B(0)
    },
    \label{eq:stack_slide_retained_power}
\end{equation}
and the semicoherent mismatch is
\begin{equation}
    \mu_{\rm sc} = 1-R .
    \label{eq:semicoherent_mismatch}
\end{equation}

\subsubsection{Semicoherent metric}
\label{sec:semicoherent_metric}

For small offsets, the bin response is locally quadratic,
\begin{equation}
    \frac{B(d)}{B(0)}
    =
    1-\kappa d^2 + {\cal O}(d^4).
\end{equation}
For the unaveraged rectangular-window response, $\kappa=\pi^2/3$.  For
$B_{\rm avg}$, $\kappa$ is the curvature of
$B_{\rm avg}(d)/B_{\rm avg}(0)$ at $d=0$.  Substitution into
Eq.~\eqref{eq:stack_slide_retained_power} gives a local metric on the
semicoherent track parameters,
\begin{equation}
    \mu_{\rm sc}
    \simeq
    g^{\rm sc}_{ab}\Delta\theta^a\Delta\theta^b,
\end{equation}
where
\begin{equation}
    g^{\rm sc}_{ab}
    =
    \kappa
    \frac{
        \sum_{i\in{\cal I}(\theta)}
        w_i
        \left[
            \Delta f_{{\rm bin},i}^{-1}
            \frac{\partial f_i}{\partial\theta^a}
        \right]
        \left[
            \Delta f_{{\rm bin},i}^{-1}
            \frac{\partial f_i}{\partial\theta^b}
        \right]
    }{
        \sum_{i\in{\cal I}(\theta)} w_i
    } .
    \label{eq:semicoherent_metric_tensor}
\end{equation}
This metric is the power-stack analogue of the coherent mismatch metric.  It
keeps the same geometric interpretation---nearby templates are close if their
parameter offset has small metric norm---but the norm measures loss of
stack-slide power rather than loss of coherent matched-filter SNR.

In practice the derivatives $\partial f_i/\partial\ln M_c$ are evaluated by
finite differences of TaylorF2 tracks, while
$\partial f_i/\partial t_{20}=-\dot f_i$ is obtained from the same frequency
evolution.  The sums in Eq.~\eqref{eq:semicoherent_metric_tensor} are evaluated
at the actual coherent chunk centers that contribute to the track, rather than
by a continuous time average.

\subsubsection{Template-count estimate}
\label{sec:semicoherent_count_estimate}

The metric volume of the searched track space is
\begin{equation}
    V_{\rm sc}
    =
    \int_\Theta
    \sqrt{\det g^{\rm sc}(\theta)}
    \,d\ln M_c\,dt_{20}.
    \label{eq:semicoherent_metric_volume}
\end{equation}
For each chirp mass, the allowed range of $t_{20}$ is
\begin{equation}
    -T_{\rm band}(M_c) \leq t_{20} \leq T_{\rm obs},
    \label{eq:t20_range}
\end{equation}
where $T_{\rm band}(M_c)$ is the TaylorF2 time required to chirp across the
analysis band.  This convention includes signals that crossed $20\,{\rm Hz}$
before the start of the observation, provided that some part of the track lies
inside the observed data.

For a target maximum stack-slide mismatch $\mu_{\max}$, the metric volume can
be converted into a template-count estimate by dividing by a covering cell
area.  In two dimensions,
\begin{equation}
    A_{\rm square} = 2\mu_{\max},
    \qquad
    A_{\rm hex} = \frac{3\sqrt{3}}{2}\mu_{\max},
\end{equation}
for square and hexagonal coverings, respectively.  Thus
\begin{equation}
    N_{\rm square}\simeq\frac{V_{\rm sc}}{A_{\rm square}},
    \qquad
    N_{\rm hex}\simeq\frac{V_{\rm sc}}{A_{\rm hex}}.
    \label{eq:semicoherent_template_counts}
\end{equation}
These counts are intended to set the order of magnitude of the stack-slide
problem and to guide the computational-cost estimates.

The calculation does not, by itself, produce a final search bank.  In
particular, it does not implement the binary-tree placement, metric-reuse
rules, metric-failure recovery, boundary padding, or edge-effect handling used
in MANIFOLD-style production banks.  Those details are important for a search
whose efficiency must be validated at the boundaries of the target parameter
space.  The present metric should therefore be viewed as the statistic-specific
ingredient needed by a bank generator, not as a complete bank generator.  A
production bank could be built by adapting MANIFOLD to use
Eq.~\eqref{eq:semicoherent_metric_tensor}, or by validating a PyCBC brute-force
or stochastic bank directly against the retained-power criterion in
Eq.~\eqref{eq:stack_slide_retained_power}.
